\phantomsection\label{fig:vol01:western-zhou:zhao-southern-routes}
\begin{center}
\begin{tikzpicture}[x=.88cm,y=.88cm,font=\small]
  \fill[paper] (0,0) rectangle (11.2,6.0);

  \draw[source!65!timeline,line width=1pt]
    (6.7,.35) .. controls (7.4,1.4) and (7.0,2.6) .. (8.2,3.55)
    .. controls (9.1,4.25) and (9.0,5.0) .. (10.5,5.65);
  \node[text=source!80!ink,right] at (8.35,2.2) {汉水};
  \draw[source!42,dashed] (3.1,5.15) .. controls (4.2,4.65) and (5.1,4.7) .. (5.9,4.2);
  \node[text=source!80!ink,above] at (4.5,4.85) {秦岭南麓方向};

  \fill[dynasty!18] (.7,3.75) -- (2.7,3.55) -- (3.0,4.85) -- (1.0,5.15) -- cycle;
  \node[align=center] at (1.85,4.35) {宗周\\关中};
  \fill[timeline!20] (3.95,2.7) -- (5.85,2.55) -- (6.15,3.85) -- (4.25,4.1) -- cycle;
  \node[align=center] at (5.05,3.3) {南阳—汉水\\北部走廊};
  \fill[source!14] (7.35,.65) -- (10.35,.75) -- (10.55,2.65) -- (8.1,2.95) -- cycle;
  \node[align=center] at (9.1,1.75) {汉水以南\\接触区域};

  \draw[-{Stealth[length=2mm]},ultra thick,color=dynasty]
    (2.8,4.2) -- (3.78,3.65);
  \draw[-{Stealth[length=2mm]},ultra thick,color=timeline]
    (6.2,3.05) .. controls (6.9,2.55) and (7.35,2.45) .. (8.0,2.2);
  \node[text=dynasty,above,align=center] at (3.15,4.05) {王室调动\\与补给};
  \node[text=timeline,above,align=center] at (7.05,2.85) {军旅、物资与\\信息链延伸};

  \draw[-{Stealth[length=2mm]},thick,dashed,color=source]
    (8.3,2.7) .. controls (9.65,3.2) and (10.1,3.9) .. (9.55,4.6);
  \node[draw=dynasty!75,fill=paper,rounded corners=2pt,align=center,inner sep=4pt]
    at (9.35,5.0) {汉水失利的\\传统叙事所系\\地点未定};
  \node[text=source,align=center,font=\footnotesize] at (4.3,.85)
    {接触可由器物与考古旁证；\\不应还原为一条连续、确定的远征路线};
\end{tikzpicture}

\smallskip
\small\textit{图\quad 昭王南向行动与汉水风险示意：用宗周至汉水以南的交通、补给与信息距离，说明多次南向接触和后来“丧师汉水”叙事所面对的空间约束。参考《史记·周本纪》、今本《竹书纪年》、《左传·昭公十三年》以及南方青铜器、汉水流域考古；行动次数、目标、行军路径和失利地点均有争议。所有边界与路线均为约略示意。}
\end{center}
